\answersfor{ch:hsldetect}

\begin{solution}{ex:hsldetect:1}
Self-completion: $\hat x_2=(0.50+0.40)/2=0.45$, residual $0.45-0.45=0$. Spatial:
\begin{align*}
\vect{r} &= (0.50-0.20,\ 0.45-0.30,\ 0.40-0.40) = (0.30,\ 0.15,\ 0),\\
\norm{\vect{r}}^2 &= 0.09+0.0225+0 = 0.1125 .
\end{align*}
Invisible to the self-completer, plainly visible to the spatial one.
\end{solution}

\begin{solution}{ex:hsldetect:2}
\begin{align*}
1{,}050{,}895\times149 &= 156{,}583{,}355,\\
156{,}583{,}355\times4.8713\times10^{-6} &= 762.8 \approx 763 .
\end{align*}
The run logged 736, 27 fewer: about one standard deviation ($\sqrt{763}\approx27.6$).
\end{solution}

\begin{solution}{ex:hsldetect:3}
$1-(1-0.01)^5=1-0.99^5=1-0.95099=0.0490$, about 4.9\%.
\end{solution}

\begin{solution}{ex:hsldetect:4}
Objects: the interval $[-4.0,+1.5]$ points contains zero, so no difference is shown. pAUC: $[-0.0589,-0.0318]$ lies
entirely below zero, so HyperSL RM's pAUC is lower, resolved.
\end{solution}

\begin{solution}{ex:hsldetect:5}
Chance is what scores unrelated to the objects would catch. Catching fewer means the objects' scores are systematically
\emph{lower} than the background's: HyperSL's loss is smaller on the anomalies, so it reconstructs them better than
ordinary ground. That is self-completion, and it matches the AUCs below 0.5 of the 2026-09-23 run.
\end{solution}

\begin{solution}{ex:hsldetect:6}
$10^6\times438\times10^{-6}=438$\,s, 7.3 minutes. With one draw, about a quarter, $\approx110$\,s, assuming the time is
spent per draw\tagA{} (each draw runs the tokenizer, encoder and decoder once).
\end{solution}
